\scanpage{39}
We shall introduce another structure of a transcendental nature, by introducing the universal covering $\bar S_{\mathbb C}$ of $S(\mathbb C)$ at $\xi$ $\xi$ [understood: of the connected component of $\xi$ in $S(\mathbb C)$].

Let $\unsure{\gamma|\mathscr{Y}}_{G}$ be the set of homomorphisms from $\mathbb S=\prod_{\mathbb C/\mathbb R}\mathbb G_{m\,\mathbb C}$ to $G_{\mathbb R}$, corresponding to the central homomorphism \deleted{given} $i:\mathbb G_m \relbar G$ given by $i(\lambda)=i_1(\lambda)i_2(\lambda)$. \margin{\rule{0.35pt}{2.1\baselineskip}}We regard $\gamma$ as a \deleted{\unsure{analytic|analytic}} holomorphic manifold on which \deleted{$\unsure{\gamma|G}$} $G(\mathbb R)$ and a fortiori $G(\mathbb Q)$ act. This being so, we have a \uline{map}
\[
\boxed{f:\bar S_{\mathbb C}\longrightarrow\gamma_G}
\]
satisfying the conditions

\uline{12)}a. \ a) $f(\bar\xi\inserted{\text{the marked point of }\bar S_{\mathbb C}})={}\text{the homomorphism }\mathbb S\relbar G_{\mathbb R}$ defined by $i_1^{\mathbb C},i_2^{\mathbb C}$ of c).

b) $f(g\bar s)=\unsure{w_{1}|w}(g)f(\bar s)$ if $g\in\pi_1(S(\mathbb C),\xi),\ \bar s\in\bar S_{\mathbb C}$

c) $f$ is \uline{holomorphic}.

\unsure{14)|15)} (Assuming Hodge) \deleted{For every $\bar s\in\bar S_{\mathbb C}$, the largest algebraic subgroup of $G$ which over $\mathbb R$ contains $f(\bar s)(\mathbb S)\subset G_{\mathbb R}$, and which moreover contains $H$} \inserted{The composite \deleted{\unsure{of}} $\bar S_{\mathbb C}\longrightarrow\gamma_G\longrightarrow\gamma_{\widetilde G}$} (which a priori \inserted{\unsure{in view of|in view of|n\textsuperscript{o}} : 12 b)} factors through \deleted{\unsure{$\bar S_{\mathbb C}$}} $S(\mathbb C)^{\circ}$) is a \uline{constant} map.

13) If $\bar s\in\bar S_{\mathbb C}$ lies over the point $t$ of $S$ [it quite \unsure{pleases|pleases} us that $P$ as a torsor with \unsure{integrable|integrable} connection $P_S$ be defined over $S$] so that we have torsor with integrable connection \deleted{\illegible{1}} $(P_t,G_{k(t)})$ over $k(t)$, \deleted{\unsure{let}} the homomorphism $f(\bar s)_{\mathbb C}:\mathbb S_{\mathbb C}\longrightarrow(G_{k(t)})_{\mathbb C}\simeq G_{\mathbb C}$ [defining a bigrading\inserted{ (hence filtration)} on the functor $T_{\mathbb C}$, $T$ forgetful functor on $G$-modules] is \uline{compatible} with the torsor $P_t$ under $G_{k(t)}$ and its $C$-trivialization $p(\bar s)$ [$p:\bar S_{\mathbb C}\longrightarrow(P_S)_{\mathbb C}$ $\unsure{\in|\epsilon}$ \unsure{here is|is} the $S(\mathbb C)^{\circ}$-morphism characterized by the condition that $p(\bar\xi)=\alpha$].
% archive sheet number: 38
